\chapter{Podsumowanie}

Stworzenie aplikacji \emph{Grocery Scanner} wymagało ogromnego nakładu czasowego. Proces jej powstawania był pełen różnych koncepcji i pomysłów na rozbudowę aplikacji. Wiele z~nich często podlegało zmianom, lecz efekt końcowy pozwala stwierdzić, że aplikacja spełnia wyrażone we wstępie założenia. Możliwe jest jednak poszerzanie jej o dodatkowe funkcjonalności. Przykładowymi planami na rozwój aplikacji są:

\begin{itemize}
    \item \emph{możliwość logowania i rejestracji za pomocą konta w takich serwisach, jak Google i~Facebook}, co~usprawni zaimplementowany już proces, a w konsekwencji pozwoli na powiększenie liczby użytkowników korzystających z aplikacji,
    \item \emph{system zgłaszania profili użytkowników}, który pozwoli na ograniczenie szkodliwej działalności w aplikacji poprzez ograniczanie lub blokowanie dostępu,
    \item \emph{rozszerzenie listy preferencji żywieniowych}, która ułatwi korzystanie z aplikacji większej liczbie użytkowników,
    \item \emph{pobieranie informacji o produktach z innych źródeł}, co pozwoli na wzrost jakości danych, a także w dużym stopniu ograniczy przypadki, w których aplikacja zwraca informację o braku danych skanowanego produktu.
\end{itemize}

\noindent
Warto zaznaczyć, że funkcjonalności zawarte w powyższej liście nie zostały zaimplementowane z uwagi na fakt, że wykraczały poza podstawowe założenia co do funkcjonalności aplikacji opisanych we wstępie pracy.
%  w małym stopniu przybliżały do zrealizowania podstawowych założeń aplikacji zawartych we wstępie, co było głównym celem.

Proces tworzenia niniejszej pracy przede wszystkim pozwolił na pogłębienie wiedzy. W~tej kwestii szczególną uwagę należy zwrócić uwagę na aspekt „techniczny”. Tworzenie aplikacji \emph{Grocery Scanner} pozwoliło na poznanie nieużywanych wcześniej przez autora pracy bibliotek i~narzędzi, jak również umożliwiło poszerzenie wiedzy na temat poznanych wcześniej rozwiązań. Sam proces tworzenia aplikacji na urządzenia mobilne okazał się być niezwykle pouczający, zarówno w kontekście poszerzania wiedzy technicznej, jak i rozwoju kariery zawodowej.